\chapter{Building a music player}
\label{ch:musicplayer}

\begin{aside}
Library browser, playlist, controls. The music itself plays
via an external library (libmpv, libVLC, ffmpeg).
\end{aside}

\section{Layout}

Three panels: library tree (left), playlist (centre),
now-playing card (right). Plus a transport bar at the
bottom.

\section{Library}

The library is a tree of artist > album > track. Build
the tree from a metadata cache; rescan on demand.

\section{Playlist}

A reorderable list. Drag-and-drop or up/down with shift.

\section{Playback}

A background thread (or external process) handles audio.
\maya{} sends commands (play, pause, seek, next) and
receives state (position, duration, current track).

\section{Now playing}

Track title, artist, album art (if image protocol),
progress bar, elapsed/remaining time.

\section{Visualisation}

A live waveform or spectrum analyser as a sparkline.
Cheap eye candy.

\section{Controls}

\begin{itemize}[leftmargin=*]
  \item Space: play/pause.
  \item Left/Right: seek.
  \item J/K: previous/next.
  \item Up/Down: volume.
\end{itemize}

\section{Shuffle and repeat}

Modes toggled with R and S. Icons in the transport bar.

\section{Library search}

\code{/} focuses a search box. Fuzzy-match against title,
artist, album.

\section{Recap}

\begin{recap}
\begin{itemize}[leftmargin=*]
  \item Library tree + playlist + now-playing.
  \item External audio engine; \maya{} drives it.
  \item Transport bar with controls.
  \item Optional visualisation.
  \item Search across metadata.
\end{itemize}
\end{recap}

\section*{Workshop}

\begin{enumerate}[leftmargin=*]
  \item Build the layout shell with mock data.
  \item Wire an audio engine (libmpv).
  \item Add visualisation.
\end{enumerate}
